\documentclass{article}
\usepackage{amsmath,amsfonts,graphicx}
\usepackage[a4paper, left=2cm, right=2cm, top=2cm, bottom=2cm]{geometry}
\usepackage{hyperref}
\usepackage{listings}
\usepackage{xcolor}
\usepackage{graphicx}
\setlength{\parskip}{10pt}
\setlength{\parindent}{0pt}

% Modern styling
\setlength{\parskip}{10pt}
\setlength{\parindent}{0pt}
\linespread{0.5}


\definecolor{codegray}{gray}{0.95}
\lstset{
	backgroundcolor=\color{codegray},
	basicstyle=\ttfamily\small,
	frame=single,
	breaklines=true
}


\def\img_dir{"../reports/images"}


\begin{document}
	% Titleplage
%	\pagestyle{empty}
	\begin{center}
		\LARGE
		\textbf{Report: root-to-shoot ratio} \vspace{0.5cm}
		\hrule
	\end{center}
	\tableofcontents
	\bigskip  
	\bigskip  
	\hrule
	
	\section{Summary}
	
		\textbf{Data Exploration}
		\begin{itemize}
			\item Root mass slightly more right-skewed then sprout mass
			\item Sprout-root-mass-ratio shows stronger peak around 0.5 - 0.7
			\item High correlation between sprout- and root-mass  ($r = 0.76$)
		\end{itemize}
		
		\textbf{Marginal Distribution}
		\begin{itemize}
			\item Gamma fits data better then Normal, support with Kullback-Leibler-Divergence
		\end{itemize}
		
		\textbf{Ratio Modeling}
		\begin{itemize}
			\item The theoretical model doesn't match the empirical distribution of the ratios
			\item Supports statistical dependency of root- and sprout-mass
		\end{itemize}
		
		\textbf{Interpretation}
		\begin{itemize}
			\item The growth of root and sprout might be statistical dependent (with statisitcal power)
			\item Since the theoretical model under estimates the peak, one could argue that this ratio range might be a optima for breeding
			\item Since I don't have much information about the data, I cannot interpret much about the biology
		\end{itemize}
		
	\newpage
	
	\section{Data Exploration}

		\subsection*{Data Distribution}
	
		
			\begin{center}
				\includegraphics[scale=0.5]{"\img_dir/01_histograms.png"}
			\end{center}
			
			\textbf{Masses}
			\begin{itemize}
				\item The masses of sprouts and roots tend to be slightly right skewed (rechtsschief)  
				\item $Weights \in \mathbb{R}_+$   
				\item Just a normal distribution might now be appropriate, because its highly likely it will extend into the negatives and doesn't catch the skewness  
				\item A Gamma might fits the marginal distribution better, because: $X \in \mathbb{R}_{+}$
			\end{itemize}
		
			\textbf{Ratio}
			\begin{itemize}
				\item The skweness seem to not be present here   
				\item $Ratio \in [0, 1]$    
				\item Intuitivly (and by their definition) a beta might fit best here  
			\end{itemize}
		
		\subsection*{Data Correlation}
			\begin{center}
				\includegraphics[scale=0.5]{"\img_dir//01_correlation.png"}
			\end{center}
			
			\begin{itemize}
				\item Hight correlation between the masses (=0.76)
				\item Implying a statistical dependency between both parameters !!!
			\end{itemize}	

	
	
	\section{Marginal Distribution}
		
		\subsection*{Gamma fits data better then a Normal}
			\begin{center}
				\includegraphics[scale=0.5]{"\img_dir/02_est_distr.png"}		
			\end{center}
			
			\begin{itemize}
				\item Just visually I would argue that the Gamma fits  $X_{sprout}$ and $X_{root}$ better then a Normal  
				\item Especially, since a Gamma can model the skweness of data
				\item One can clearly see, that the normal extends into the negatives
			\end{itemize}
			
		\subsection*{Further Evaluation with Kullback-Leibler-Divergence}
		
			\textbf{Sprout}	
			\begin{center}
				\includegraphics[scale=0.5]{"\img_dir//02_discr_sprout.png"}	
			\end{center}
				
			\begin{center}
				\begin{tabular}{ccc}
					    & Gamma & Normal \\
					KLD: & 0.01999 & 0.05624
				\end{tabular}
			\end{center}	
			\begin{itemize}
				\item Gamma has a much lower KLD 
			\end{itemize}
			\newpage
			
			\textbf{Root}	
			\begin{center}
				\includegraphics[scale=0.5]{"\img_dir//02_discr_root.png"}	
			\end{center}
			
			\begin{center}
				\begin{tabular}{ccc}
					& Gamma & Normal \\
					KLD: & 0.02861 & 0.06414
				\end{tabular}
			\end{center}	
		
			\begin{itemize}
				\item Also here, Gamma has a much lower KLD 
			\end{itemize}
		
		\textbf{Estimated parameters from Gamma (sprout and root)}
		\begin{verbatim}
			Params.       Value         
			----------------------------
			a_sprout      3.4580
			b_sprout      0.2310
			a_root        2.0912
			b_root        0.2266
		\end{verbatim}		

		\begin{itemize}
			\item One can notice that:  $\beta_{root} \approx \beta_{sprout}$
			\item This supports the theoretical distrubtion for the ratio transformation later
			\item But also that:  $\alpha_{root} < \alpha_{sprout}$, implying that the skewness of the root distribution is stronger !
		\end{itemize}
		\newpage
	
		
	\section{Ratio-Modelling}
	
		\subsection*{Root-To-Sproot Ratio}
			
			\textbf{Theoretically}
			
			\begin{tabular}{ll}
				$\quad X_1 \sim Ga(x_1| \alpha_1, \beta)$ & Shown in previous section \\
				$\quad X_2 \sim Ga(x_2| \alpha_2, \beta)$ & Shown in previous section \\
				$\quad Y = \frac{X_1}{X_1 + X_2}$ & Transformation \\
				$\quad Y \sim Be(y|\alpha_1, \alpha_2)$ & Under the assumption that $\beta_{root} = \beta_{sprout} = \beta$ and $X_1 \perp\!\!\!\perp X_2 $  (independent). \\ 
				& This was proven in our lectures.  
			\end{tabular}
			
			- I can theoretically model the ratio with just the $\alpha$'s from the Marginal Gamma's
	
		\begin{center}
			\includegraphics[scale=0.6]{"\img_dir/03_ratio_beta_modelling.png"}
		\end{center}
	
		\begin{itemize}
			\item One can clearly see that the theoretically derived beta distribution doesn't fit the empirical ratios
		\end{itemize}
		
	\subsection*{Interpretation}
	\begin{itemize}
		\item The poor fit quality of the estimated beta distribution  is an indicator that both masses are statistically dependent and further supports the correlation 
		\item This validates, that the correlation breaks the independency assumption of the theoretical beta
		\item One could interpret that this peak is the result of stronger dependency (also correlation) of those seedlings masses
		\item Unfortunately, I don't have to much information about the data, especially I don't know if certain seedlings perform better than other (for breeding purpose)
		\item I could argue that this ratio range (0.5 - 0.7) might be a optimal target for breeding selection (with statistical power)
	\end{itemize}
	
	
\end{document}